% =====================================================================
% Resumen (IMRyD, unas 300 palabras) y abstract. Redactado el 06/11/2026
% (sesión R4). Las cifras proceden del capítulo 5.
% =====================================================================
\chapter*{Resumen}
\addcontentsline{toc}{chapter}{Resumen}

\textbf{Introducción.} La previsión del día siguiente de la demanda eléctrica
horaria sirve de base a la casación del mercado diario. La comparación entre
modelos clásicos y de aprendizaje automático suele hacerse sin igualdad de
condiciones y solo con el error puntual. Este trabajo las compara con un
protocolo equitativo para la demanda peninsular española y evalúa también la
calibración de los intervalos de predicción.

\textbf{Métodos.} Se construye un conjunto horario de 2017 a 2024 con la demanda
de ESIOS, la temperatura de reanálisis ponderada por población y el calendario
laboral. Se comparan un ingenuo estacional, un SARIMAX con términos de Fourier,
LightGBM y la red N-HiTS, con la misma información y el mismo presupuesto de
búsqueda. 2023 sirve para elegir y 2024, con reajuste mensual y una sola
ejecución, para evaluar. La precisión se mide con MAE, RMSE, MAPE, MASE y el
contraste de Diebold--Mariano; los intervalos del 90\,\%, propios y con
inferencia conforme adaptativa (ACI), con cobertura, puntuación de intervalo y
CRPS. La temperatura de reanálisis equivale a una previsión meteorológica
perfecta.

\textbf{Resultados.} En 2024 el MAE es de 573,6~MW en LightGBM, 665,6 en N-HiTS,
1089,7 en el SARIMAX y 1217,9 en el ingenuo. LightGBM y N-HiTS superan de forma
significativa al SARIMAX y al ingenuo, y LightGBM supera a N-HiTS; el SARIMAX no
se distingue del ingenuo con la regla fijada de antemano. Los intervalos propios
de LightGBM y N-HiTS infracubren (76,1 y 78,7\,\% para un 90\,\% nominal) y los
de los modelos clásicos sobrecubren por poco. Con ACI los cuatro quedan entre el
88 y el 92\,\%, pero ACI no mejora una ventana de residuos recientes. Ningún
régimen cambia de forma distinguible el modelo preferible, y los modelos
clásicos fallan sobre todo en festivos.

\textbf{Discusión.} En este problema el aprendizaje automático mejora al SARIMAX,
y la calibración depende más de cómo se construye el intervalo que del modelo.
Los resultados son de un sistema, un año y cuatro modelos, y una cota superior
por la previsión meteorológica perfecta.

\medskip
\noindent\textbf{Palabras clave:} previsión de la demanda eléctrica; series
temporales; SARIMAX; aprendizaje automático; intervalos de predicción.

\clearpage
\begin{otherlanguage}{english}
\chapter*{Abstract}
\addcontentsline{toc}{chapter}{Abstract}

\textbf{Introduction.} Day-ahead forecasts of hourly electricity demand underpin
the day-ahead market. Comparisons between classical and machine-learning models
are often made on unequal terms and with point error alone. This thesis compares
them under a fair protocol for peninsular Spanish demand and also evaluates the
calibration of prediction intervals.

\textbf{Methods.} An hourly dataset for 2017--2024 combines ESIOS demand,
population-weighted reanalysis temperature and the working-day calendar. A
seasonal naïve benchmark, a SARIMAX with Fourier terms, LightGBM and the N-HiTS
network are compared with the same information and search budget. 2023 is used to
choose and 2024, with monthly refitting and a single run, to evaluate. Accuracy is
measured with MAE, RMSE, MAPE, MASE and the Diebold--Mariano test; 90\,\%
intervals, each model's own and those built with adaptive conformal inference
(ACI), with coverage, interval score and CRPS. Reanalysis temperature amounts to
a perfect weather forecast.

\textbf{Results.} In 2024 the MAE is 573.6~MW for LightGBM, 665.6 for N-HiTS,
1089.7 for SARIMAX and 1217.9 for the naïve model. LightGBM and N-HiTS
significantly outperform SARIMAX and the naïve model, and LightGBM outperforms
N-HiTS; SARIMAX cannot be distinguished from the naïve model under the rule fixed
in advance. The own intervals of LightGBM and N-HiTS undercover (76.1 and
78.7\,\% for a nominal 90\,\%) and those of the classical models slightly
overcover. With ACI all four fall between 88 and 92\,\%, but ACI does not improve
on a window of recent residuals. No regime changes the preferred model in a
distinguishable way, and the classical models fail mainly on public holidays.

\textbf{Discussion.} In this problem machine learning improves on SARIMAX, and
calibration depends more on how the interval is built than on the model. The
results refer to one system, one year and four models, and are an upper bound
because of the perfect weather forecast.

\medskip
\noindent\textbf{Keywords:} electricity demand forecasting; time series; SARIMAX;
machine learning; prediction intervals.
\end{otherlanguage}
